%\section{System prototype}
\section{PID interoperability with NDN}
\label{pid-poc}

% introduction of what's happening here

%Short intro about it
%-Show our model
%\subsection{PID interoperability}
In this section, we will propose a design, which achieves \gls{pid} interoperability with the \gls{ndn} namespace and makes it feasible to add future \gls{pid} types to answer our research question in section \ref{introduction-research-question}. Our design is based on related work done by Karakannas \cite{icn-bd}, by avoiding the \gls{pid} to \gls{ndn} translation on client side \cite{icn-bd}. The pattern matching method was based on related work by Mousa for identifying different \gls{pid} type schemas \cite{ndn-app-aware}. The research done by Olschanowsky et. al. was used for deriving \gls{ndn} names from metadata \cite{ndn-man}.
We have reimplemented the principles of \gls{naas4pid} in order to demonstrate the \gls{pid} interoperability extensibility feature. We came up with our own solution, as the \gls{naas4pid} was nor available or public to extend it.
We combined the aforementioned related work and extended it with an extendable interoperability feature, we attempt to make the translation transparent to the user and support multiple \gls{pid} types. The proposed design for our solution is illustrated in figure \ref{fig:sdc_model}.
Our design adheres to the following aforementioned principles, which will be discussed in more detail in this section.
 
\begin{itemize}
    \item{Translation is transparent to the user.}
    \item{Support for multiple \gls{pid} types.}
    \item{Extensible with future \gls{pid} types with different naming schemas.}
\end{itemize}

\begin{figure}[H]
\centering
\includegraphics[width=\textwidth]{Images/PIDtoNDN11.png}
\caption{NDN virtual function based 
planning for achieving \gls{pid} interoperability}
\label{fig:sdc_model}
\end{figure}

Our proposed design consists of the following components, each with their own functionality; the \gls{pid} server, the \gls{pid} to \gls{ndn} gateway and the client. The general idea is that a user enters a \gls{pid} of the object that the user want to retrieve at the client and gets back the requested object as shown in figure \ref{fig:sdc_model}. The retrieval of an object depends if the object is already published in the \gls{ndn} or not, which is further described in sections \ref{client} and \ref{gw}.
%If it is published in NDN, then the client retrieves the object from \gls{ndn}. If it is not published in NDN, the object is retrieved from the \gls{pid} server. This is further described in section \ref{gw}.  

% integrate with section 3, planning concept
\subsection{Proof of concept} 
This section discusses the implementation of our design in a proof of concept.
In our proof of concept, the components mentioned in section \ref{pid-poc} are conceptually part of the \gls{ndn} we have setup in our high-level network design, which will be discussed in section \ref{planning-architecture}.

\subsubsection{PID server}
The \gls{pid} server is maintained by the \gls{pid} provider, which identifies objects of a particular \gls{pid} type. The \gls{pid} types we cover in our proof of concept are highlighted in section \ref{pid-types}. For our proof of concept we have set up a Handle \gls{pid} server with Cordra software \cite{cor} on our lab environment server 'nimes' to store and identify our data. We got allotted the Handle prefix \texttt{20.5000.481} by the Handle registry. We use this prefix for object identification. The Handles stored on our Handle \gls{pid} server are resolved by the Handle System (\texttt{http://hdl.handle.net}). In addition to this, we also used the resolver of the national library of the Netherlands for resolving URN's as well as PANGAEA, which resolves \glspl{doi}. 

\subsubsection{Client}\label{client}
The role of the client is to transparently make the user either retrieve the requested object from the corresponding \gls{pid} provider or from the \gls{ndn}. This will be further described in section \ref{gw}. The client receives the \gls{pid} based on the user's input, which can be any kind of \gls{pid} type. The client then opens a socket and sends the user's request to the gateway, which corresponds with step 1 and 2 in figure \ref{fig:sdc_model}. After the gateway does the translation, it send back either a translated \gls{ndn} name from the \gls{pid} or a link to the \gls{pid} server to the client. This depends whether the object is published in the \gls{ndn} or not. If it is, the client receives the \gls{ndn} name. If not, it receives the link to the \gls{pid} server for retrieving the object. This corresponds to step 6a in figure \ref{fig:sdc_model}.

If an \gls{ndn} name is sent back to the client, the client will request the object from the NDN\footnote{\url{https://github.com/AquaL1te/rp2/blob/master/Scripts/ndn_client.py}}. This is shown as step 7a and step 8 till 10 in figure \ref{fig:sdc_model}. For object retrieval in \gls{ndn} we used the tool \texttt{ndncatchunks} part of the \texttt{ndn-tools} software, which makes use of the libraries of the \gls{cxx} application \cite{ndn-tools}. 
If the object is not in \gls{ndn} and a \gls{pid} link is sent back, the client requests the object by its \gls{pid} link by sending the request to the \gls{pid} server\footnote{\url{https://github.com/AquaL1te/rp2/blob/master/Scripts/pid_client.py}}. This corresponds to step 7c in figure \ref{fig:sdc_model}. Object retrieval is done at client side, otherwise the gateway has to retrieve the object first before sending it to the client. 
%As in the case of retrieving the object from the \gls{pid} server to publish it in NDN, 
The client would have to wait for the gateway till it has retrieved and cached the object in \gls{ndn}. 
In the case that the object is already published in NDN, the gateway server only concerns itself with translation and sending back the \gls{ndn} name to the client. This eliminates unnecessary load on the gateway.

Retrieving an object either from the \gls{pid} server or from \gls{ndn} should happen transparently for the user. It is possible to accomplish this by combining the code for retrieving the object from the \gls{pid} server with the code we have used to retrieve the object from \gls{ndn} at the client side. Furthermore, a conditional statement needs to be added at the gateway to check if the object is already published in \gls{ndn}. This is further described in section \ref{gw}.
%can be done with the \texttt{ndnping} tool \cite{ndn-tools} for example as described in section \ref{gw}.

As a result, if transparency is implemented the user might not even be aware of an \gls{ndn}. The user only specifies a \gls{pid} as input for the client, without specifying the link of the web resolver that resolves the \gls{pid}. The user gets redirected automatically by the client to either the \gls{pid} server or \gls{ndn}. No further user input is needed as everything is taken care of by the gateway, which will also be discussed in section \ref{gw}.
%Unlike previous designs, which require user input after the client has translated the \gls{pid} to \gls{ndn} name \cite{ndn-app-aware}. 

\subsubsection{Gateway}\label{gw}
The gateway used in our design follows the principles of Karakannas by avoid doing the \gls{pid} to \gls{ndn} translation on client side. In the case of doing the translation on client side, the client software needs to be updated every time a new \gls{pid} type is introduced. For our proof of concept, we designed a translation server called the \gls{pid} to \gls{ndn} gateway. The \gls{pid} to \gls{ndn} gateway implements the translation of different \gls{pid} types and sends the translated name back to the client. Furthermore, we identify \gls{pid} types based on pattern matching as described by Mousa \cite{ndn-app-aware}. We have also looked at how the \gls{n2t} resolver deals with different \gls{pid} types. The \gls{n2t} resolves different \gls{pid} types by stating the \gls{pid} type that needs to be implemented along with the pattern of the PIDs' schema \cite{n2t}.

%Translation
The gateway is responsible for translating a \gls{pid} to \gls{ndn} name and checks if the requested object is already published in \gls{ndn}. 
The first responsibility is translating the \gls{pid} it receives from the client to an \gls{ndn} name, which can be of any \gls{pid} type. In our proof of concept we have implemented the Handle \gls{pid} type schema of the Handle \gls{pid} server we setup. In addition to this, we have also implemented the \gls{urn} \gls{pid} type schema of the national library of the Netherlands, as well as the \gls{doi} type schema of PANGAEA. The gateway receives a \gls{pid} from the client, without the link of the corresponding web resolver of the \gls{pid} type. Based on pattern matching of the \gls{pid} type schema\footnote{\url{https://github.com/AquaL1te/rp2/blob/master/Scripts/pid_server.py\#L58-L62}}, the gateway detects what kind of \gls{pid} type it has to deal with. Then, the associated function is called to translate the \gls{pid} to \gls{ndn} name\footnote{\url{https://github.com/AquaL1te/rp2/blob/master/Scripts/pid_server.py\#L17-L37}} and appends the corresponding link of the web resolver of the \gls{pid} type it receives. 
%Pattern matching is done based on the patterns of 
The patterns of most standardized \gls{pid} type schemas are maintained in the ePIC \gls{dtr} \cite{dtr} and can be used for implementation in the gateway we propose. 

%Naming:
Furthermore, the second responsibility of the gateway is to check if the object is already published in NDN.
%This can be done with a ping to the object with \texttt{ndnping}, part of the \texttt{ndn-tools} software \cite{ndn-tools} used in our proof of concept. To use this, a \texttt{ndnpingserver} has to be setup for the objects with the \gls{ndn} testing software we have utilized. The presence of the object can also be detected with \texttt{ndncatchunks} by catching the error exception, which is thrown if the object is not published by \texttt{ndnputchunks}.
If the object is available in NDN, the gateway sends the translated \gls{ndn} name back to the client. The client then retrieves the object from \gls{ndn}. This is shown in figure \ref{fig:seq_ndn}, where the Handle \gls{pid} type is used as an example. If the object is not available in NDN, the gateway sends back the \gls{pid} link to the client and caches the object in \gls{ndn}. 
%The client then retrieves the object by its \gls{pid} link from the \gls{pid} server as shown in figure \ref{fig:seq_pid}. 
The \gls{pid} link that is sent back to the client contains the \gls{pid} and the link to the corresponding \gls{pid} web resolver. This happens before the gateway retrieves the object to cache it in NDN, otherwise the client needs to wait for this as discussed in section \ref{client}. For our proof of concept we use the tool \texttt{ndnputchunks}, which is part of the \texttt{ndn-tools} software for caching objects in the \gls{ndn} \cite{ndn-tools}. Transparency can be achieved, if the code we used in our setup to retrieve the object from \gls{ndn} is combined with the code we used to retrieve the object from the \gls{pid} server. This is done by including a conditional statement to check if the object is already published in \gls{ndn}. This can be done with \texttt{ndnping} for example, which is also part of the \texttt{ndn-tools} testing software that we have utilized. To use \texttt{ndnping}, a \texttt{ndnpingserver} has to be setup for the objects. Another option is to catch the error, which is thrown by \texttt{ndncatchunks} if the object is not published by \texttt{ndnputchunks}. 
To translate the matched \gls{pid} type to an \gls{ndn} name, one has to take into account the schema of the \gls{pid} type and the hierarchical way it has to be divided in. In our proof of concept a prefix is added before the \gls{pid} name, for deriving an \gls{ndn} name from each discussed \gls{pid} type. Such as \texttt{/ndn/handle} for Handle objects, \texttt{/ndn/doi} for \gls{doi} objects and \texttt{/ndn/urn} for \gls{urn} objects. This complies with the related work we have discussed in section \ref{introduction-related-work}. The delimiters of the \gls{pid} types are replaced by a slash \texttt{("/")}. In PANGAEA, specific columns and parameters can be requested to retrieve a particular part from an object\footnote{\url{https://doi.pangaea.de/10.1594/PANGAEA.842227&columns=1,,2,3&filterParameterValue=Station,TARA_100}}. The columns and parameters can also be translated to an \gls{ndn} name\footnote{\url{/ndn/doi/10.1594/PANGAEA.842227/attrib+ndn+1,2,3+Station,TARA_100}} \cite{ndn-app-aware}. This shows that not only the delimiters are replaced to translate it into an \gls{ndn} name. This is done by hierarchically dividing it in \texttt{/attrib+ndn} and appending the columns and parameters with a plus \texttt{("+")} (where the requested columns and parameters are separated by a comma \texttt{(","))}. 
Afterwards, the requested data chunk is stored in \gls{ndn} with the parameter and column attributes. The next time someone wants to retrieve that part, it will be already available. Another solution would be to request the whole object and hierarchically divide the attributes in an \gls{ndn} name by some kind of rules. This leads to having all chunks of a object available in NDN, so the next time a part of the object can be requested.
Furthermore, the web resolver link is also stripped after translation.

The web resolver link is not used for deriving the \gls{ndn} name.
By excluding the web resolver link from the \gls{ndn} name, duplications will not occur in \gls{ndn} as the \gls{ndn} name is only derived from the \gls{pid}. A \gls{pid} always translates to the same name in \gls{ndn} this way. 
If the \gls{pid} object is moved to another web resolver, only the link of the web resolver has to be updated in the gateway. 

In listing \ref{lst:hdl_ndn}, the translation of a Handle \gls{pid} to an \gls{ndn} name is shown, the Handle name is hierarchically divided into its \gls{pid} type, authority and sub-authority.
\vspace{1em}
\begin{lstlisting}[frame=single,gobble=0,basicstyle=\scriptsize\ttfamily,caption={Handle \gls{pid} to \gls{ndn} name}\label{lst:hdl_ndn}]
<user>@consumer-1:~/python-ndn$ python3 server_pid.py
Waiting for client
PID from connected user: 20.500.481/sub-auth/object1
PID type: Handle
NDN name from Handle: /ndn/handle/20.500.481/sub-auth/object1
\end{lstlisting}

The translation of a \gls{urn} \gls{pid} to an \gls{ndn} name is shown in listing \ref{lst:urn_ndn} for the \gls{anp} collection maintained by the national library of the Netherlands as discussed in section \ref{urn-1}. The national library of the Netherlands chose to assign a \gls{pid} based on the year when the object has been published. 
This way, objects in \gls{ndn} can be hierarchically divided in year, months or days for example.
\vspace{1em}
\begin{lstlisting}[frame=single,gobble=0,basicstyle=\scriptsize\ttfamily,caption={URN \gls{pid} to \gls{ndn} name}\label{lst:urn_ndn}]
<user>@consumer-1:~/python-ndn$ python3 server_pid.py
Waiting for client
PID from connected user: anp:1938:10:01:2:mpeg21
PID type: URN
NDN name from URN: /ndn/urn/anp/1938/10/01/2/mpeg21
\end{lstlisting}

Metadata can be used if missing gaps need to be filled in for dividing the \gls{pid} in the \gls{ndn} hierarchy, as described by Olschanowsky et. al. \cite{ndn-clim} in section \ref{introduction-related-work}. 
Filling in missing gaps highly depends on which way the metadata is served by the providers of the different \gls{pid} types as discussed in section \ref{pid-types}. Thus, if metadata is used, a parser has to be implemented in the gateway. This could be a XML, JSON or an other kind of parser depending on the \gls{pid} provider. In our proof of concept we have implemented a XML parser for URN's and a JSON parser for Handles.

There is no address exhaustion problem in \gls{ndn} as the \gls{ndn} namespace is unbounded \cite{ndn-nspace}. But worth mentioning is to keep in mind that using long \gls{ndn} names degrades performance with many interests as described by Yuan et al. \cite{yuan2012scalable} in section \ref{introduction-related-work}.

\begin{figure}[H]
%\flushleft
    \centering
%    \makebox[0pt]{%
    \includegraphics[width=\textwidth]{Images/pid_seq5.png}
    %}
    \caption{Handle \gls{pid} request\label{fig:seq_pid}}
\end{figure}

\begin{figure}[H]
\includegraphics[scale=0.75]{Images/ndn_req.png}
\caption{Handle \gls{ndn} request}
\label{fig:seq_ndn}
\end{figure}


% merge with results
\subsection{Results}\label{results-pid}
The outcome of implementing our design in a proof of concept shows that our principles can be adhered. Since the \gls{naas4pid} solution was not available to extend or improve it, we came up with our own solution and made it publicly available. Making the translation and object retrieval transparent to the user is possible. Users should not have to concern themselves whether to retrieve the object through \gls{ndn} or the \gls{pid} server. This is due to gateway's responsibility for \gls{pid} to \gls{ndn} translation and the object retrieval, which is taken care of by the client. 
%This can be achieved by combining the code we used for retrieving the object from either the \gls{pid} server or \gls{ndn}. 
Translation is achieved by first recognizing the \gls{pid} type based on pattern matching and then hierarchically divide the \gls{pid} to an \gls{ndn} name. Support for multiple \gls{pid} types is also achieved by adding the schema of the \gls{pid} types at the gateway, which makes it also easily extensible to support future \gls{pid} types. By adding \gls{pid} types at the gateway, we overcome the hurdle of updating the client software with the schema of newly introduced \gls{pid} types each time when a new \gls{pid} type is introduced.




%\subsection{Deploying an \gls{ndn} network (McCabe and TOSCA)}
%\label{planning-deploying}
% refer to sections where architecture and design is defined
%Now that the network analysis, design and architecture are defined, a deployment strategy is needed. The high-level design (figure \ref{fig:high-level-network-design}) needs to become deployable with a scalable method. Scalable in this context means that a single deployment strategy can be used for different cloud providers. Furthermore, as defined in the scope (section \ref{introduction-scope}) if the infrastructure scales out, the effort for managing a larger infrastructure should be equal. As described in section \ref{overview-tosca}, \gls{tosca} is a standard to describe the complete life cycle of an infrastructure. Having a single set of template descriptions for deployment benefits portability and reproducibility of an infrastructure on different cloud providers.

%\begin{figure}[H]
%\centering
%\includegraphics[width=\columnwidth]{Images/tosca-diagram.png}
%\caption{TOSCA diagram.}
%\label{fig:tosca-diagram}
%\end{figure}

%In figure \ref{fig:tosca-diagram}, a \gls{tosca} diagram is illustrated. This diagram represents an abstract template description of the \gls{tosca} relationships, in which the grey rectangular boxes are the core scalability factors. As described in section \ref{overview-tosca}, \gls{tosca} consists out of several types; nodes, relationships and interfaces. The scaling properties are highlighted in the rectangular areas. The left area, highlighted as 'scaling in/out resources' contains a dependency chain of several virtual \gls{ndn} functions. This dependency chain is also depicted numerically. Before a pod can be deployed on Kubernetes (step 2 to 5), a VM needs to exist (step 1). This is described by the 'dependsOn' relationship. Furthermore, with the requirements defined in section \ref{planning-requirements}, input constraints are described. These constraints are used by the orchestrator to make sure that the \gls{ndn} infrastructure has sufficient resources available to operate. Once a VM is deployed, the dependency for Kubernetes is satisfied, thus Kubernetes can then be setup (step 2). Kubernetes can then deploy pods by the use of interfaces (step 3). These interfaces feed the containers with environment variables such as the gateway, a list of routes, the transport protocol for NDN, the \gls{ndn} strategies and on which Kubernetes node this pod should run. The environment variables are given to the interface via the \gls{tosca} inputs. These environment variables are then used by scripts that run inside the pods to setup \gls{ndn}. Several constraints are set for these environment variables such as which valid transport protocols can be used for NDN, which \gls{ndn} strategies are valid and which nodes are available. These constraints are defined with e.g. 'valid\_values' or 'greater\_than' definitions. These constraints help to guide the orchestrator to verify the inputs that are given for the template description. As illustrated in the second gray area 'scaling in/out the application', several pods can be instantiated (step 5a, 5b and 5c) from the image (step 4). These pods enable the virtual \gls{ndn} functions as described in section \ref{planning-architecture}. These pods establish the \gls{ndn} network and therefore are connected via the 'connectsTo' relationship. This network expands over to other Kubernetes nodes in the cluster by the use of the Kubernetes built-in overlay network.

%\subsection{PID interoperability architecture design}
%The discussed architecture for our proposed design in section \ref{planning-architecture} is shown in figure \ref{fig:sdc_model}. This section will explain more about the different components used in our proposed design.

%\begin{figure}[H]
%\centering
%\includegraphics[width=\textwidth]{Image%s/PIDtoNDN10.png}
%\caption{NDN virtual function based 
%planning for achieving \gls{pid} interoperability}
%\label{fig:sdc_model}
%\end{figure}


%\subsubsection{PID server}
%The \gls{pid} server is maintained by the \gls{pid} provider, which identifies objects of a particular \gls{pid} type. The \gls{pid} types we cover in our proof of concept are highlighted in section \ref{pid-types}. For our proof of concept we have set up our own Handle \gls{pid} server with Cordra software \cite{cor}. We got allotted the Handle prefix \texttt{20.5000.481} by the Handle registry. We use this prefix for object identification. In addition to this, we also used the resolver of the national library of the Netherlands for resolving URN's as well as PANGAEA, which resolves DOI's. 

%\subsubsection{Client}\label{client}
%The role of the client is to transparently make the user either retrieve the requested object from the corresponding \gls{pid} provider or from the \gls{ndn} network. This will be further described in this section \ref{gw}. The client receives the \gls{pid} based on the user's input, which can be any kind of \gls{pid} type. The client then opens a socket and sends the user's request to the gateway, which corresponds with step 1 and 2 in figure \ref{fig:sdc_model}. After the gateway does the translation, it send back either a translated \gls{ndn} name from the \gls{pid} or a link to the \gls{pid} server to the client. This depends whether the object is published in the \gls{ndn} network or not. If it is, the client receives the \gls{ndn} name. If not, it receives the link to the \gls{pid} server for retrieving the object. This corresponds to step 6a in figure \ref{fig:sdc_model}.

%If an \gls{ndn} name is sent back to the client, the client will request the object from the \gls{ndn} network\footnote{\url{https://github.com/AquaL1te/rp2/blob/master/Scripts/ndn_client.py}}. This is shown as step 7a and step 8 till 10 in figure \ref{fig:sdc_model}. For object retrieval in \gls{ndn} we used the tool \texttt{ndncatchunks} part of the \texttt{ndn-tools} software \cite{ndn-tools}, which makes use of the libraries of the \gls{cxx} application \cite{ndn-tools}. 
%If the object is not in \gls{ndn} and a \gls{pid} link is sent back, the client requests the object by its \gls{pid} link by sending the request to the \gls{pid} server\footnote{\url{https://github.com/AquaL1te/rp2/blob/master/Scripts/pid_client.py}}. This corresponds to step 7c in figure \ref{fig:sdc_model}. Object retrieval is done at client side, otherwise the gateway has to retrieve the object first before sending it to the client. 
%As in the case of retrieving the object from the \gls{pid} server to publish it in NDN, 
%The client would have to wait for the gateway till it has retrieved and cached the object in \gls{ndn}. 
%In the case that the object is already published in NDN, the gateway server only concerns itself with translation and sending back the \gls{ndn} name to the client. This eliminates unnecessary load on the gateway.

%Retrieving an object either from the \gls{pid} server or from \gls{ndn} should happen transparently for the user. It is possible to accomplish this by combining the code for retrieving the object from the \gls{pid} server with the code we have used to retrieve the object from \gls{ndn} at the client side. Furthermore, a conditional statement needs to be added at the gateway to check if the object is already published in \gls{ndn}. This is further described in section \ref{gw}.
%can be done with the \texttt{ndnping} tool \cite{ndn-tools} for example as described in section \ref{gw}.

%As a result, if transparancy is implemented the user might not even be aware of an \gls{ndn} network. The user only specifies a \gls{pid} as input for the client, without specifying the link of the web resolver that resolves the \gls{pid}. The user gets redirected automatically by the client to either the \gls{pid} server or \gls{ndn} network. No further user input is needed as everything is taken care of by the gateway, which will be discussed below in section \ref{gw}.
%Unlike previous designs, which require user input after the client has translated the \gls{pid} to \gls{ndn} name \cite{ndn-app-aware}. 

%\subsubsection{Gateway}\label{gw}
%The gateway used in our design follows the principles of Karakannas by avoid doing the \gls{pid} to \gls{ndn} translation on client side. In the case of doing the translation on client side, the client software needs to be updated every time a new \gls{pid} type is introduced. For our proof of concept, we designed a translation server called the \gls{pid} to \gls{ndn} gateway. The \gls{pid} to \gls{ndn} gateway implements the translation of different \gls{pid} types and sends the translated name back to the client. Furthermore, we identify \gls{pid} types based on pattern matching as described by Mousa \cite{ndn-app-aware}. We have also explored the \gls{n2t} resolver, which resolves different \gls{pid} types by stating the \gls{pid} type that needs to be implemented along with the pattern of the PIDs' schema \cite{n2t}.

%Translation
%The gateway is responsible for translating a \gls{pid} to \gls{ndn} name and checks if the requested object is already published in \gls{ndn}. 
%The first responsibility is translating the \gls{pid} it receives from the client to an \gls{ndn} name, which can be of any \gls{pid} type. In our proof of concept we have implemented the Handle \gls{pid} type schema of the Handle \gls{pid} server we setup. In addition to this, we have also implemented the \gls{urn} \gls{pid} type schema of the national library of the Netherlands, as well as the \gls{doi} type schema of PANGAEA. The gateway receives a \gls{pid} from the client, without the link of the corresponding web resolver of the \gls{pid} type. Based on pattern matching of the \gls{pid} type schema\footnote{\url{https://github.com/AquaL1te/rp2/blob/master/Scripts/pid_server.py\#L58-L62}}, the gateway detects what kind of \gls{pid} type it has to deal with. Then, the associated function is called to translate the \gls{pid} to \gls{ndn} name\footnote{\url{https://github.com/AquaL1te/rp2/blob/master/Scripts/pid_server.py\#L17-L37}} and appends the corresponding link of the web resolver of the \gls{pid} type it receives. 
%Pattern matching is done based on the patterns of 
%The patterns of most standardized \gls{pid} type schemas are maintained in the ePIC \gls{dtr} \cite{dtr} and can be used for implementation in the gateway we propose. 

%Naming:
%Furthermore, the second responsibility of the gateway is to check if the object is already published in NDN.
%This can be done with a ping to the object with \texttt{ndnping}, part of the \texttt{ndn-tools} software \cite{ndn-tools} used in our proof of concept. To use this, a \texttt{ndnpingserver} has to be setup for the objects with the \gls{ndn} testing software we have utilized. The presence of the object can also be detected with \texttt{ndncatchunks} by catching the error exception, which is thrown if the object is not published by \texttt{ndnputchunks}.
%If the object is available in NDN, the gateway sends the translated \gls{ndn} name back to the client. The client then retrieves the object from \gls{ndn}. This is shown in figure \ref{fig:seq_ndn}, where the Handle \gls{pid} type is used as an example. If the object is not available in NDN, the gateway sends back the \gls{pid} link to the client and caches the object in \gls{ndn}. 
%The client then retrieves the object by its \gls{pid} link from the \gls{pid} server as shown in figure \ref{fig:seq_pid}. 
%The \gls{pid} link that is sent back to the client contains the \gls{pid} and the link to the corresponding \gls{pid} web resolver. This happens before the gateway retrieves the object to cache it in NDN, otherwise the client needs to wait for this as discussed in section \ref{client}. For our proof of concept we use the tool \texttt{ndnputchunks}, which is part of the \texttt{ndn-tools} software for caching objects in the \gls{ndn} network \cite{ndn-tools}. Transparency can be achieved, if the code we used in our setup to retrieve the object from \gls{ndn} is combined with the code we used to retrieve the object from the \gls{pid} server. This is done by including a conditional statement to check if the object is already published in \gls{ndn}. This can be done with \texttt{ndnping} for example, which is also part of the \texttt{ndn-tools} testing software that we have utilized. To use \texttt{ndnping}, a \texttt{ndnpingserver} has to be setup for the objects. Another option is to catch the error, which is thrown by \texttt{ndncatchunks} if the object is not published by \texttt{ndnputchunks}. 
%To translate the matched \gls{pid} type to an \gls{ndn} name, one has to take into account the schema of the \gls{pid} type and the hierarchical way it has to be divided in. In our proof of concept a prefix is added before the \gls{pid} name, for deriving an \gls{ndn} name from each discussed \gls{pid} type. Such as \texttt{/ndn/handle} for Handle objects, \texttt{/ndn/doi} for \gls{doi} objects and \texttt{/ndn/urn} for \gls{urn} objects. This complies with the related work we have discussed. The delimiters of the \gls{pid} types are replaced by a slash \texttt{("/")}. In PANGAEA, specific columns and parameters can be requested to retrieve a particular part from an object\footnote{\url{https://doi.pangaea.de/10.1594/PANGAEA.842227&columns=1,,2,3&filterParameterValue=Station,TARA_100}}. The columns and parameters can also be translated to an \gls{ndn} name\footnote{\url{/ndn/doi/10.1594/PANGAEA.842227/attrib+ndn+1,2,3+Station,TARA_100}} \cite{ndn-app-aware}. This shows that not only the delimiters are replaced to translate it into an \gls{ndn} name. This is done by hierarchically dividing it in \texttt{/attrib+ndn} and appending the columns and parameters with a plus \texttt{("+")} (where the requested columns and parameters are separated by a comma \texttt{(","))}. Furthermore, the web resolver link is also stripped after translation.

%The web resolver link is not used for deriving the \gls{ndn} name.
%By excluding the web resolver link from the \gls{ndn} name, duplications will not occur in \gls{ndn} as the \gls{ndn} name is only derived from the \gls{pid}. A \gls{pid} always translates to the same name in \gls{ndn} this way. 
%If the \gls{pid} object is moved to another web resolver, only the link of the web resolver has to be updated in the gateway. 

%In listing \ref{lst:hdl_ndn}, the translation of a Handle \gls{pid} to an \gls{ndn} name is shown, the Handle name is hierarchically divided into its \gls{pid} type, authority and sub-authority.
%\vspace{1em}
%\begin{lstlisting}[frame=single,gobble=0,basicstyle=\scriptsize\ttfamily,caption={Handle \gls{pid} to \gls{ndn} name}\label{lst:hdl_ndn}]
%<user>@consumer-1:~/python-ndn$ python3 server_pid.py
%Waiting for client
%PID from connected user: 20.500.481/sub-auth/object1
%PID type: Handle
%NDN name from Handle: /ndn/handle/20.500.481/sub-auth/object1
%\end{lstlisting}

%The translation of a \gls{urn} \gls{pid} to an \gls{ndn} name is shown in listing \ref{lst:urn_ndn}. The national library of the Netherlands chose to assign a \gls{pid} based on the year when the object has been published. 
%This way, objects in \gls{ndn} can be hierarchically divided in year, months or days for example.
%\vspace{1em}
%\begin{lstlisting}[frame=single,gobble=0,basicstyle=\scriptsize\ttfamily,caption={URN \gls{pid} to \gls{ndn} name}\label{lst:urn_ndn}]
%<user>@consumer-1:~/python-ndn$ python3 server_pid.py
%Waiting for client
%PID from connected user: anp:1938:10:01:2:mpeg21
%PID type: URN
%NDN name from URN: /ndn/urn/anp/1938/10/01/2/mpeg21
%\end{lstlisting}

%Metadata can be used if missing gaps need to be filled in for dividing the \gls{pid} in the \gls{ndn} hierarchy, as described by Olschanowsky et. al. \cite{ndn-clim} in section \ref{introduction-related-work}. 
%Filling in missing gaps highly depends on which way the metadata is served by the providers of the different \gls{pid} types as discussed in section \ref{pid-types}. Thus, if metadata is used, a parser has to be implemented in the gateway. This could be a XML, JSON or an other kind of parser depending on the \gls{pid} provider. In our proof of concept we have implemented a XML parser for URN's and a JSON parser for Handles.

%There is no address exhaustion problem in \gls{ndn} as the \gls{ndn} namespace is unbounded \cite{ndn-nspace}. But worth mentioning is to keep in mind that using long \gls{ndn} names degrades performance with many interests as described by Yuan et al. \cite{yuan2012scalable} in section \ref{introduction-related-work}.

%\begin{figure}[H]
%\flushleft
%    \centering
%    \makebox[0pt]{%
%    \includegraphics[width=\textwidth]{Images/pid_seq5.png}
    %}
%    \caption{Handle \gls{pid} request\label{fig:seq_pid}}
%\end{figure}

%\begin{figure}[H]
%\includegraphics[scale=0.75]{Images/ndn_req.png}
%\caption{Handle \gls{ndn} request}
%\label{fig:seq_ndn}
%\end{figure}

% merge with pid setup
%\subsection{Proof of concept}
%\label{planning-poc}
%With the methodology defined, in which scalability and performance requirements are met and a method for deployment is described, a proof of concept was used to test the methodology. The orchestrators mentioned in section \ref{overview-tosca} are still in a prototype phase. Therefore, in our proof of concept we deployed the VMs and Kubernetes nodes manually. In practice the life cycle of the Kubernetes pods are managed by a \gls{tosca} orchestrator. Without having a \gls{tosca} orchestrator available, steps 2 through 5 in figure \ref{fig:tosca-diagram} were be carried out by Kubernetes exclusively. This was done by defining the configuration properties\footnote{\url{https://github.com/AquaL1te/rp2/blob/master/Kubernetes/expanded-cluster.yml}} of the pods manually. These properties include the \gls{ndn} function name, e.g. router, producer or consumer. And also includes the routes (NDN prefixes) and the associated \gls{ndn} face with the transport protocol to use (TCP or UDP). These parameters were then inserted into the \gls{ndn} \gls{fib} by the scripts that were executed inside the pod\footnote{\url{https://github.com/AquaL1te/rp2/blob/master/Docker/producer/docker-entrypoint.sh}}. The \gls{ndn} strategies were also configured by these scripts. Furthermore, if it is not defined where a pod should be running, Kubernetes will make this decision itself, based on the known resources in the Kubernetes cluster. If for example a Kubernetes node has more memory to spare than other nodes, then Kubernetes will likely decide to spawn the pod there. This Kubernetes node could potentially run in a cloud provider, located in another geographical area. Since the purpose is to provide data distribution through the use of NDN, locality becomes a key factor. Therefore, a pod is specifically assigned to a Kubernetes node in order to provide in-network caching in a specific geographical area.

% merge with pid results
%\subsection{Results}
%As a result, the \gls{ndn} infrastructure life cycle can be managed from Kubernetes. Our proof of concept lacks a \gls{tosca} orchestrator. Therefore, scaling in or out resources to other cloud provider is not demonstrated. However, scaling in or out the \gls{ndn} application is demonstrated. This method allows to reconfigure an \gls{ndn} infrastructure in NVF-style by interacting with Kubernetes as the orchestrator. Therefore, the YAML configuration of Kubernetes acts as the \gls{tosca} template description and Kubernetes, which executes this configuration, acts as the orchestrator. In practice the Kubernetes configuration would be generated based on the \gls{tosca} template descriptions and e.g. \gls{drip} would act as the orchestrator. In effect, the \gls{ndn} containers are spawned as NVFs which provides scalable management of the \gls{ndn} network.

%The outcome of implementing our design %for \gls{pid} interoperability 
%in a proof of concept shows that our principles can be adhered. Making the translation transparent to the user is possible, which means that the user may not be aware of receiving the object from the \gls{pid} server or \gls{ndn}. This is due to gateway's responsibility for \gls{pid} to \gls{ndn} translation and the object retrieval, which is taken care of by the client. This can be achieved by combining the code we used for retrieving the object from either the \gls{pid} server or \gls{ndn}. Translation is achieved by first recognizing the \gls{pid} type based on pattern matching and then hierarchically divide the \gls{pid} to an \gls{ndn} name. Support for multiple \gls{pid} types is also achieved by adding the schema of the \gls{pid} types at the gateway, which makes it also easily extensible to support future \gls{pid} types. By adding \gls{pid} types at the gateway, we overcome the hurdle of updating the client software with the schema of newly introduced \gls{pid} types each time when a new \gls{pid} type is introduced.